\section{Summary}\label{sec:offline-summary}

This chapter addresses offline retail CBDC payments: payments conducted
while both payer and payee are disconnected from the network, and settled
only when tokens are eventually deposited. The central difficulty in this
setting is that a payee has no way to verify, in real time, that the token
they are receiving has not already been spent elsewhere. Rather than rely on
prevention or detection alone, we combine the two: a secure element makes
double spending expensive and impractical to mount at scale, while an
efficient auditing mechanism ensures that, should an attack succeed, the
responsible party can still be identified and held accountable.

We construct this scheme from anonymous credentials, blind signatures, and
zero-knowledge proofs of knowledge, which together hide the identities of
payers and payees from the central bank and the intermediaries and prevent
withdrawals from being linked to deposits, and from a verifiable encryption
scheme that lets an authorised auditor efficiently de-anonymise double
spenders without compromising the privacy of honest users. Since these
primitives are too costly for off-the-shelf secure elements, we split the
computation of a payment between a lightweight operation performed by the
secure element and a heavier one offloaded to the user's mobile device, so
that the secure element performs work comparable to producing an ordinary
digital signature.

We implement this split using BBS+ signatures for the anonymous credential
issued by the registration authority, and Pointcheval-Sanders blind
signatures for the tokens issued by the central bank, and demonstrate the
viability of the resulting protocol on commodity smart cards and mobile
phones. Withdrawal and the payer side of a payment both complete in under
half a second, and even a payment chain of 32 hops verifies in under a
second on the payee's side, with the smart card itself authorising a
payment in under $400\,\mathrm{ms}$. Auditing adds no overhead for the
payer and a modest 10-20\% overhead for the payee.

Deploying the scheme in practice raises questions we have only discussed
informally in this chapter. Revoking users without undermining their
privacy, and without exposing payees who are offline at the moment of
revocation, is a genuine challenge in this setting; the time-to-live and
transaction-value restrictions we discuss go some way towards mitigating
the resulting window of exposure, but a more thorough treatment of
revocation, closer to the level of formality we give the rest of the
scheme, would be valuable future work. Likewise, our security analysis
assumes an idealised secure element; a concrete evaluation of the residual
risk once realistic hardware constraints and attacker budgets are taken
into account would help translate our guarantees into deployment guidance.

We focus throughout on transferability rather than divisibility, following
the same trade-off made by related e-cash schemes discussed in
Section~\ref{sec:offline-related}; extending the scheme to support splitting
and merging tokens, and to support aggregated deposits as more than the
two-token case we sketch in Appendix~\ref{sec:aggregate-deposits}, are both
natural directions for future work.
